\section{智能的模态分化：资源、时间与确定性的约束三角}

\subsection{约束三角的三条边与张力}

残差落在哪个位置，决定了吸收它所需要的资源结构、时间尺度与确定性水平，而这三种要求并不总是同向的。智能不会只沿着一条直线进化。一种常见想象是：智能越强，就越通用、越快速、越可靠、越低成本，最后收敛成某种统一的全能系统，它能实时理解世界，瞬间生成方案，稳定控制现实，准确预测后果，并以极低资源消耗完成所有任务。这个想象把智能看成一种单轴能力：能力越大，一切维度同步改善。但真实世界中的智能不是在真空中增长，而是在约束中运行，任何智能系统都必须处理资源、时间和确定性之间的张力，它需要资源来计算、感知、存储、冷却和执行，需要时间来接收信息、推理、验证、反馈和修正，还需要确定性来保证结果可靠、行为稳定、风险可控、责任可追踪；三者可以被技术不断改善，却很难同时无代价推向极值。\key{单轴想象的关键错误，是以为约束会随能力一起消失。}

\subhead{约束三角的三条边}

这不是一个数学意义上的绝对不可能定理，而是一个物理和系统工程意义上的约束判断。更好的算法、更高效的芯片、更强能源、更优传感器、更成熟的控制架构，当然都能移动边界。但移动边界不等于取消边界。技术可以把约束三角推大，却不能让速度、确定性和低成本同时无限提高而不产生新的代价。

三条边各自有明确的物理含义。所谓资源，是智能运行的物理账本：计算需要能量，模型需要存储，传感需要设备，推理需要时间片，执行需要功率，控制需要通信和反馈，高可靠系统还需要冗余、备份、审计和维护。即便一个智能系统看起来像纯信息结构，它也必须运行在芯片、电力、网络、冷却和制造体系上。智能越要覆盖更大系统、更高频反馈、更高后果任务，它越难保持轻盈。

数据中心把智能的物理底座显示出来。IEA在2026年4月与10月的报告分别检查效率、需求与供电条件：单位任务效率可以改善，同时总电力需求仍受使用增长推动。\cite{iea2026questions,iea2026powering} 比较时要区分全部数据中心、AI相关设施和单项任务，不能把早期预测当成2026年的实际结果。

把总量分解成需求与单位能耗两项，两个方向的证据就可以放进同一个式子，

\begin{equation}
E=De,\qquad \frac{1}{E}\frac{dE}{dt}=\frac{1}{D}\frac{dD}{dt}+\frac{1}{e}\frac{de}{dt}
\end{equation}

$E$ 是给定周期内的总用电量，单位为太瓦时；$D$ 是同一周期内被满足的算力需求，单位为折算后的算力单位；$e$ 是单位算力的能耗，单位为太瓦时每算力单位。左侧是第一式的量分解，右侧由两边取对数后求导得到，三项都是无量纲的年增长率。这个式子的含义是：“效率在改进”与“这一条边仍在移动”可以同时为真，只要需求增长率大于单位能耗下降的幅度，总用电量的增长率就仍然为正。把它用在正文的两组数据上，八年之间算力需求增至 6.5 倍而用电量仅增至 1.06 倍，单位算力能耗因此降至原来的约六分之一，平均每年下降约两成，两项才能同时成立。适用条件是分解本身可加：当算力需求与能耗的统计口径不一致（例如训练、推理与冷却分别计量），或需求发生结构性替换时，$e$ 的含义随口径变化，恒等式仍然成立，但两侧不再能被独立解释。这个式子说明的是账本结构，不是对能耗趋势的预测。

所谓时间，是智能与世界交换信息的过程。智能不能在信息尚未到达之前知道所有后果，也不能在过程尚未展开之前完成真实验证。快速响应往往意味着在不完整信息下行动，长期验证往往意味着放慢速度，复杂反馈往往意味着必须等待系统返回结果。越追求即时决策，越容易依赖局部视野；越追求全局鲁棒，越需要更多观察、校验和反馈周期。

验证时间由不同部分组成：证明、仿真、试验准备、真实暴露、随访与审定。更好的工具可以加速其中一些部分；老化、长期反应与足量观察则有各自条件。项目需要逐项说明哪些等待在产生证据，哪些等待只是流程拥堵。

所谓确定性，是系统对错误、风险和后果的控制能力。低风险文本生成可以容忍模糊、噪声和再生成，高风险控制系统不能这样。自动驾驶、核电站、电网、金融清算、手术机器人、航空控制、药物审批，都不能只看平均表现，而必须关心长尾、边界、责任、可解释、可审计和可恢复。确定性不是一句“模型更聪明”就能获得，它必须通过限制、冗余、验证和制度结构来购买。

FDA在2026年9月的官方介绍中称，已有超过1600项AI医疗器械获得美国市场授权。\cite{fda2026aidevices} 清单定期更新，并明确说明识别方式与覆盖限制。\cite{fda2026ailist} 每项授权对应具体产品与用途；数量表明一批受限定功能已经进入应用，不等于一种通用医疗资格。
\begin{figure}[ht]
\centering
% >>>>>>>>>> figs/fig-modality-triangle.tex >>>>>>>>>>
% 图 4-1  约束三角：资源、时间、确定性
% 三个顶点是三笔刚性支出，任意两类智能只能同时靠近其中两个，
% 第三条边必然被牺牲。三类智能的位置因此不是偏好选择，而是被三条边夹出来的。
%   探索型靠近「资源—时间」边：便宜、快，但结论不确定。
%   控制型靠近「时间—确定性」边：确定、及时，但要付出冗余与功率。
%   验证型靠近「资源—确定性」边：确定、便宜，但慢。
% 三条边各写一个可测量的量：功率密度、验证时延、冗余度。
\begin{tikzpicture}[x=1cm, y=1cm, font=\small,
  ax/.style={ink, line width=0.8pt},
  vert/.style={circle, fill=ink, inner sep=1.8pt},
  pt/.style={circle, inner sep=2.4pt},
  ee/.style={font=\scriptsize, text=slate, align=center}
]
\useasboundingbox (-0.60,0.85) rectangle (12.75,-7.55);

% 顶点：上=资源，左下=时间，右下=确定性
\def\vR{6.10} \def\vTop{-0.55} \def\vBot{-5.35} \def\vHalf{4.35}
\coordinate (A) at (\vR,\vTop);                 % 资源
\coordinate (B) at ({\vR-\vHalf},\vBot);        % 时间
\coordinate (C) at ({\vR+\vHalf},\vBot);        % 确定性
\fill[slate!8] (A) -- (B) -- (C) -- cycle;
\draw[ax] (A) -- (B) -- (C) -- cycle;
% 中线（辅助读位置，不参与论证）
\draw[slate!30, line width=0.4pt, dashed] ({\vR},{\vBot}) -- ({\vR},{\vBot+2.90});
\draw[slate!30, line width=0.4pt, dashed] (B) -- ({\vR+\vHalf/2},{\vTop/2+\vBot/2});
\draw[slate!30, line width=0.4pt, dashed] (C) -- ({\vR-\vHalf/2},{\vTop/2+\vBot/2});

\node[vert] at (A) {};
\node[vert] at (B) {};
\node[vert] at (C) {};
\node[font=\small\heiti, text=ink, anchor=south] at (\vR,{\vTop+0.16}) {资源};
\node[font=\small\heiti, text=ink, anchor=north east] at ({\vR-\vHalf-0.08},{\vBot-0.12}) {时间};
\node[font=\small\heiti, text=ink, anchor=north west] at ({\vR+\vHalf+0.08},{\vBot-0.12}) {确定性};

% ---------- 三条边上的可测量 ----------
\node[ee, anchor=east] at ({\vR-\vHalf/2-0.14},{\vTop/2+\vBot/2}) {功率密度\\（资源买时间）};
\node[ee, anchor=west] at ({\vR+\vHalf/2+0.14},{\vTop/2+\vBot/2}) {冗余度\\（资源买确定性）};
\node[ee, anchor=north] at ({\vR},{\vBot-0.42}) {验证时延\\（确定性买时间）};

% ---------- 三类智能的位置 ----------
\node[pt, fill=forest] at (5.10,-3.55) {};
\node[pt, fill=deepblue] at (7.05,-3.55) {};
\node[pt, fill=terra] at (6.10,-4.60) {};
\node[font=\scriptsize, text=forest!85, anchor=east, align=right] at (4.98,-3.55)
      {探索型\\便宜、快、结论不可靠};
\node[font=\scriptsize, text=deepblue!85, anchor=west, align=left] at (7.17,-3.55)
      {控制型\\及时、确定、代价是冗余};
\node[font=\scriptsize, text=terra!90, anchor=south] at (6.10,-4.72)
      {验证型：便宜、确定、但慢};

% ---------- 张力：往一个顶点靠，另两条边同时被拉长 ----------
\draw[ink!55, line width=0.6pt, -{Stealth[length=1.7mm]}]
      (6.10,-4.12) -- (6.10,-2.34);
\node[font=\scriptsize, text=ink, anchor=north, align=center] at (6.10,-1.22)
      {越靠近某个顶点\\另两条边同时变长};

% ---------- 右下角：这条约束的读法 ----------
\node[font=\scriptsize, text=slate, anchor=north west, align=left]
      at ({-0.35},{-6.72})
      {三角不是价值判断，是账本：想要确定的结论又要快，就得把资源堆上去；\\
       资源有限时，唯一能选的是把不确定性放在哪一环。};
\end{tikzpicture}

% <<<<<<<<<< figs/fig-modality-triangle.tex <<<<<<<<<<

\caption{约束三角。资源、时间与确定性是三笔刚性支出，任意一类智能最多能同时靠近其中两个顶点，第三条边必然被牺牲。三类智能各自的位置因此不是偏好选择，而是被三条边夹出来的结果：靠近一个顶点，就要放弃另外两个方向。}
\label{fig:modality-triangle}
\end{figure}


\subhead{约束三角的张力}

于是，一个智能系统越想快，就越可能牺牲验证深度；越想稳，就越需要牺牲速度和自由度；越想低成本，就越必须限制任务范围、责任范围和运行频率。这不是说快、稳、省不能同时进步，而是说三者之间存在结构性张力，每一次技术进步都能局部降低张力，但不会把张力彻底抹平，这就是智能的\term{约束三角}。这种张力在资源分配上是可以记账的
\begin{equation}
B_{\mathrm{E}}+B_{\mathrm{C}}+B_{\mathrm{V}}=B
\end{equation}
式中 $B$ 是系统在给定周期内可支配的总投入，$B_{\mathrm{E}}$、$B_{\mathrm{C}}$、$B_{\mathrm{V}}$ 分别是探索、控制、验证三类工作在同一账本上的占用；四项必须先折算到可加的数量口径（能量、算力、人力、周期时间），因此它是一条记账恒等式，而不是最优解公式。它给出的约束很直接：只要 $B$ 不变，任何一项的增加都要从其余两项中扣除，这正是“越快、越稳、越省三者不能同时无代价改善”的写法。它的成立条件是可折算性，当某些投入无法折算进同一账本，例如制度合法性、人的责任承担和长期信任，恒等式只覆盖可折算的部分，账本之外的代价仍然存在。

在这个三角中，全能智能不是默认的稳定形态。它当然可以作为极限想象存在，但一旦进入真实系统，它就必须付费：为了快，要支付信息不完备的风险；为了稳，要支付验证和冗余成本；为了省，要放弃某些能力边界；越接近“全能”，系统越要同时承担高资源消耗、高反馈复杂度、高验证负担和高责任压力，全能不是免费上升，而是多维压力的集中。因此，真实世界中的智能更可能发生模态分化。所谓\term{模态分化}，不只是智能被分配到不同任务，而是不同智能承担不同的系统功能：探索型智能展开可能性，控制型智能封装风险，验证型智能显影错误；前者扩张候选空间，中者收束行动边界，后者暴露隐藏代价。三者不是同一种智能的强弱差异，而是现实系统中三种不同的工作。

\subsection{探索型、控制型、验证型三类智能与生态接口}

\subhead{探索型智能的功能}

第一类是探索型智能，它的核心功能是展开可能性。探索型智能生成文本、图像、代码、假设、方案、分子结构、设计草案和策略候选。它的价值不在于每一次输出都确定为真，而在于低成本产生大量候选，让系统看到原本没有想到的路径。它适合低后果、低执行成本、可回滚、可筛选的比特域，也适合科学发现和工程设计的早期搜索阶段。

探索型智能必须保持流动性。它需要发散、联想、类比、偶然和试探，不能被过度锁死在确定性要求中，否则会失去生成新可能性的能力。幻觉、噪声和不稳定输出，在高风险场景中是严重缺陷，但在探索场景中，它们也部分来自低约束生成系统的结构特征。探索不是验证，候选不是事实，想象不是现实。探索型智能的合理位置，是把可能性推出来，而不是替现实确认它们。

从系统角度看，探索型智能的功能不是稳定行动，而是扩张候选空间。它把原本隐含的组合、路径、表达和设计显露出来，让可能性从少变多，让候选从稀缺变丰富。它降低的是搜索黏滞，但也制造新的下游压力：候选越多，筛选越重；表达越便宜，可信表达越稀缺；方案越容易生成，现实验证越昂贵。所以探索型智能不能直接成为高风险系统的终点，它可以提出方案，但不能单独批准方案，可以生成候选，但不能直接承担高后果执行，可以扩大搜索空间，但不能替代验证系统。把探索型智能直接接入高风险现实，就会把低确定性的生成噪声转化为现实风险。探索型智能的主场是可能性，不是责任终点。

\subhead{控制型智能的功能}

第二类是控制型智能，它的核心功能是封装风险。控制型智能不是为了生成更多可能性，而是为了让系统在噪声、扰动、延迟和故障中保持可控，自动驾驶、机器人控制、工业产线、电网调度、航空航天、核设施、金融基础设施、医疗设备，都需要这类智能。它的首要品质不是多，而是稳，不是新奇，而是可靠，不是表达能力，而是边界管理。

控制型智能必须牺牲一部分自由。它不能随意发散，不能临场编造，不能把未验证策略直接作用于高后果系统。它需要限定运行域、冗余传感、安全阈值、异常处理、日志记录、人工接管、合规审计和责任链条。它的强大不在于“什么都敢做”，而在于清楚知道什么不能做、何时降级、何时停止、何时交还控制权。从系统角度看，控制型智能不是扩大选择，而是收束行动，它要把开放世界中的风险压进可控边界，把无限情境压进有限运行域，把不可预期扰动压进冗余和安全策略。控制不是让系统变得无风险，而是让风险被隔离、缓冲、记录、降级和追责，一个成熟控制系统的价值，常常不在它表现得多聪明，而在它把多少风险封装进安全结构中。

这种封装要付出资源代价。为了可靠，它需要更多传感器；为了安全，它需要冗余；为了低延迟，它需要边缘计算和高可靠通信；为了责任，它需要记录和审计；为了长尾，它需要大量测试和验证。控制越进入原子世界和高风险系统，越无法保持生成式智能在比特域中的轻盈。它要把自由度压缩进安全边界，把智能变成稳定的闭环。控制型智能和探索型智能不是同一种东西的简单强弱差异，前者追求状态空间收束，后者追求候选空间扩张。把二者混为一谈，是 AI 时代最危险的误解之一：前者需要约束，后者需要流动；前者必须防止错误造成后果，后者容许错误被筛掉；前者面对不可随意回滚的现实，后者可以低成本重试。

\subhead{验证型智能的功能}

第三类是验证型智能，它的核心功能是显影错误。验证型智能不以生成更多答案为目标，而以暴露答案无法承受现实的部分为目标，科学推理、数学证明、形式验证、事实核查、药物证据评估、工程可靠性审查、法律合规、风险审计和长期系统评估，都需要验证型智能。它的价值不在于更快说“可以”，而在于更可靠地说“还不能”。

验证型智能必须慢下来，不是因为它低效，而是因为验证需要证据链。它需要检查来源、寻找反例、评估边界条件、测试可重复性、比较替代解释、观察长期后果。许多真实问题不能通过一次推理完成确认：药物必须经历实验和临床反馈，工程系统必须经历压力测试和故障统计，社会政策必须经历行为适应和制度后果。验证不是生成的附属品，而是现实转化的关口。从系统角度看，验证型智能的功能不是展开可能性，也不是稳定控制，而是显影错误，它把被模型忽略的边界条件显露出来，把方案中的隐含假设显露出来，把短期成功背后的长期后果显露出来，把平均表现遮蔽的长尾风险显露出来。验证型智能的任务，是让隐藏错误重新可见。

这就是验证与生成的根本差异。生成系统倾向于把世界压缩成可输出结构；验证系统则要把输出重新放回世界，让世界显影它的不足。一个生成模型可以在几秒钟内给出治疗建议，但验证系统必须问，适用于哪类患者，有什么禁忌，证据等级如何，长期副作用是什么，谁承担风险；一个模型可以生成工程方案，但验证系统必须问，极端载荷下会怎样，材料疲劳如何，维护周期如何，失败后果如何。验证型智能也有自己的资源代价，高质量验证需要数据、实验、审计、长期观测、专家判断和制度接口，它往往不像生成那样显眼，因为它的成果是减少错误、延迟冒险、阻止失败、保留边界。一个好的验证系统不一定产生耀眼输出，但它防止系统把廉价可能性误认为可靠现实。

\subhead{智能生态的接口}

三类模态到这里可以合起来看。它们都重要，但不能互相替代：一个只会展开的系统会产生大量候选，却缺乏责任；一个只会封装的系统会稳定，却缺乏开放性；一个只会显影的系统会谨慎，却可能失去行动速度。成熟智能生态不在于取消分工，而在于建立分工之间的接口。

未来不太可能只是“一个模型接管一切”。更可能出现的是多层智能系统：生成模型提出候选，仿真系统筛选路径，控制系统执行有限动作，验证系统评估安全，审计系统记录责任，人类和制度系统定义目标边界，不同模态之间互相制约、互相补足。智能越进入现实，越需要这种分工，而不是越需要一个无边界的全能单体。这种模态分化也会改变我们对“通用智能”的理解。通用能力在低黏滞比特域中非常有价值，因为任务可以快速切换，错误可以低成本回滚，输出可以被后续筛选；但进入高黏滞现实后，通用性必须被运行域、责任边界、验证标准和控制接口重新限制。越高后果的系统，越不能只依靠开放通用性，越需要把通用能力压进专门结构，使其成为可审计、可维护、可停止的工程部件。

这不是对通用智能的否定，而是对其落地形态的修正。通用智能可能成为强大的上层搜索和协调能力，但它不应被误认为可以直接替代所有底层控制和验证结构。它可以帮助提出新的路径，却不能取消物理执行的约束；可以帮助理解复杂系统，却不能取消长期反馈；可以帮助生成判断，却不能替社会承担责任。

从这个角度看，AI 时代的关键竞争不只是模型能力，而是智能生态架构。谁能更好地组织探索、控制和验证，谁就更能把智能转化为现实能力。单一模型的能力很重要，但系统接口更重要：候选如何进入验证，验证如何约束执行，执行如何反馈模型，错误如何被记录，责任如何被追踪，风险如何被隔离，目标如何被人类和制度重新定义。约束三角还提醒我们，不要把“慢”简单视为失败，某些智能必须快，某些智能必须稳，某些智能必须慢；探索慢了会错失机会，控制慢了会造成事故，验证快了则可能放过风险。真正成熟的系统不是让所有部分同频加速，而是让不同智能模态以合适速度工作，快的部分负责扩展可能，稳的部分负责守住现实，慢的部分负责确认证据和边界。

这也是文明层面的启示。未来社会不应只奖励生成速度，也不应只追求自动化覆盖率。它还必须投资验证系统、审计制度、长期实验、责任链条和人类判断。否则，探索型智能会制造候选洪水，控制型智能会被过早推入高风险现实，验证型智能会被当作拖慢创新的障碍，结果不是更智能的文明，而是更快进入堆积、失稳和信任危机的文明。

智能的模态分化不是悲观结论，而是成熟结论。它说明智能并不会因为受约束而失去价值，恰恰相反，约束让智能找到稳定形态：探索型智能让我们看到更多可能，控制型智能让可能性进入现实而不失控，验证型智能让现实改变获得可信边界，它们共同构成智能生态，而不是彼此取代。技术可以移动约束边界，但它不能让速度、确定性和低成本同时无代价达到极值。真实世界中稳定存在的，不太可能是单一全能智能，而是多种智能模态构成的生态系统。智能演化的终点不是上帝视角，而是展开可能性、封装风险、显影错误之间的分工、接口和协同。
